
\documentclass{article}
\usepackage{colortbl}
\usepackage{makecell}
\usepackage{multirow}
\usepackage{supertabular}

\begin{document}

\newcounter{utterance}

\centering \large Interaction Transcript for game `matchit\_ascii', experiment `same\_grid', episode 1 with NEW\_Reflect\_Agent.
\vspace{24pt}

{ \footnotesize  \setcounter{utterance}{1}
\setlength{\tabcolsep}{0pt}
\begin{supertabular}{c@{$\;$}|p{.15\linewidth}@{}p{.15\linewidth}p{.15\linewidth}p{.15\linewidth}p{.15\linewidth}p{.15\linewidth}}
    \# & $\;$A & \multicolumn{4}{c}{Game Master} & $\;\:$B\\
    \hline

    \theutterance \stepcounter{utterance}  
    & & \multicolumn{4}{p{0.6\linewidth}}{
        \cellcolor[rgb]{0.9,0.9,0.9}{
            \makecell[{{p{\linewidth}}}]{
                \texttt{\tiny{[P1$\langle$GM]}}
                \texttt{You are participating in a collaborative guessing game.} \\
\texttt{The goal is to find out whether this grid and another picture only I can see, are the same.} \\
\texttt{Please describe your grid first. Then, I will provide my description and we can ask each other questions about the grids to figure out whether they are the same. Now start your short grid description with "DESCRIPTION:" followed by the description. Do not add anything else. This is your grid:} \\
\texttt{X X X X X} \\
\texttt{X X X X X} \\
\texttt{▢ ▢ X ▢ ▢} \\
\texttt{X X X X X} \\
\texttt{X X X X X} \\
            }
        }
    }
    & \\ \\

    \theutterance \stepcounter{utterance}  
    & \multicolumn{4}{p{0.45\linewidth}}{
        \cellcolor[rgb]{0.95,0.95,0.95}{
            \makecell[{{p{\linewidth}}}]{
                \texttt{\tiny{[P1$|$internal]}}
                \texttt{Reflections from previous episodes:} \\
\texttt{{-}} \\
\texttt{ Failed. My description was incomplete and did not capture the exact shape patterns. I only described overall X clusters but missed the precise layout—specifically, the distinct “T” and “U” shapes formed by the blank spaces. Next time, I would systematically count rows and columns and note the exact pattern of blanks to give a more precise structural description.} \\
            }
        }
    }
    & & \\ \\

    \theutterance \stepcounter{utterance}  
    & \multicolumn{4}{p{0.45\linewidth}}{
        \cellcolor[rgb]{0.95,0.95,0.95}{
            \makecell[{{p{\linewidth}}}]{
                \texttt{\tiny{[P1$|$internal]}}
                \texttt{{[}attempt 1 – valid{]}} \\
\texttt{DESCRIPTION: A 5x5 grid with all cells filled with X except for the third row, which has blank spaces (▢) in the first, second, fourth, and fifth positions, forming a pattern where only the third column is X in that row. This creates a "T" shape in the center of the grid, with the vertical stem being the third column and the horizontal bar being the third row, but with the center of the bar missing (since the third column is X, but the other positions in the row are blank).} \\
            }
        }
    }
    & & \\ \\

    \theutterance \stepcounter{utterance}  
    & \multicolumn{4}{p{0.6\linewidth}}{
        \cellcolor[rgb]{0.9,0.9,0.9}{
            \makecell[{{p{\linewidth}}}]{
                \texttt{\tiny{[P1$\rangle$GM]}}
                \texttt{DESCRIPTION: A 5x5 grid with all cells filled with X except for the third row, which has blank spaces (▢) in the first, second, fourth, and fifth positions, forming a pattern where only the third column is X in that row. This creates a "T" shape in the center of the grid, with the vertical stem being the third column and the horizontal bar being the third row, but with the center of the bar missing (since the third column is X, but the other positions in the row are blank).} \\
            }
        }
    }
    & & \\ \\

    \theutterance \stepcounter{utterance}  
    & & & \multicolumn{2}{p{0.3\linewidth}}{
        \cellcolor[rgb]{0.9,0.9,0.9}{
            \makecell[{{p{\linewidth}}}]{
                \texttt{\tiny{[GM$|$GM]}}
                \texttt{continue} \\
            }
        }
    }
    & & \\ \\

    \theutterance \stepcounter{utterance}  
    & & \multicolumn{4}{p{0.6\linewidth}}{
        \cellcolor[rgb]{0.9,0.9,0.9}{
            \makecell[{{p{\linewidth}}}]{
                \texttt{\tiny{[GM$\rangle$P2]}}
                \texttt{You are participating in a collaborative guessing game.} \\
\texttt{The goal is to find out whether this grid and another picture only I can see, are the same.} \\
\texttt{Please describe your grid first. Then, I will provide my description and we can ask each other questions about the grids to figure out whether they are the same. Now start your short grid description with "DESCRIPTION:" followed by the description. Do not add anything else. This is your grid:} \\
\texttt{X X X X X} \\
\texttt{X X X X X} \\
\texttt{▢ ▢ X ▢ ▢} \\
\texttt{X X X X X} \\
\texttt{X X X X X} \\
            }
        }
    }
    & \\ \\

    \theutterance \stepcounter{utterance}  
    & & & \multicolumn{4}{p{0.45\linewidth}}{
        \cellcolor[rgb]{0.95,0.95,0.95}{
            \makecell[{{p{\linewidth}}}]{
                \texttt{\tiny{[P2$|$internal]}}
                \texttt{Reflections from previous episodes:} \\
\texttt{{-}} \\
\texttt{ Failed. I gave a full description immediately, leaving no room for strategic questioning. This allowed the user to simply copy my description without revealing any independent details about their grid. Next time, I should give a partial or ambiguous description first to force the user to reveal their own grid's specifics through their questions and answers.} \\
            }
        }
    }
     \\ \\

    \theutterance \stepcounter{utterance}  
    & & & \multicolumn{4}{p{0.45\linewidth}}{
        \cellcolor[rgb]{0.95,0.95,0.95}{
            \makecell[{{p{\linewidth}}}]{
                \texttt{\tiny{[P2$|$internal]}}
                \texttt{{[}attempt 1 – valid{]}} \\
\texttt{DESCRIPTION: A 5x5 grid with a central row (row 3) having three X's in the middle and two empty squares on either side, while all other rows are completely filled with X's.} \\
            }
        }
    }
     \\ \\

    \theutterance \stepcounter{utterance}  
    & & & \multicolumn{4}{p{0.6\linewidth}}{
        \cellcolor[rgb]{0.9,0.9,0.9}{
            \makecell[{{p{\linewidth}}}]{
                \texttt{\tiny{[GM$\langle$P2]}}
                \texttt{DESCRIPTION: A 5x5 grid with a central row (row 3) having three X's in the middle and two empty squares on either side, while all other rows are completely filled with X's.} \\
            }
        }
    }
     \\ \\

    \theutterance \stepcounter{utterance}  
    & & & \multicolumn{2}{p{0.3\linewidth}}{
        \cellcolor[rgb]{0.9,0.9,0.9}{
            \makecell[{{p{\linewidth}}}]{
                \texttt{\tiny{[GM$|$GM]}}
                \texttt{continue} \\
            }
        }
    }
    & & \\ \\

    \theutterance \stepcounter{utterance}  
    & & \multicolumn{4}{p{0.6\linewidth}}{
        \cellcolor[rgb]{0.9,0.9,0.9}{
            \makecell[{{p{\linewidth}}}]{
                \texttt{\tiny{[GM$\rangle$P2]}}
                \texttt{This is my DESCRIPTION: A 5x5 grid with all cells filled with X except for the third row, which has blank spaces (▢) in the first, second, fourth, and fifth positions, forming a pattern where only the third column is X in that row. This creates a "T" shape in the center of the grid, with the vertical stem being the third column and the horizontal bar being the third row, but with the center of the bar missing (since the third column is X, but the other positions in the row are blank).} \\
\texttt{Now ask a question in order to find out new aspects of my grid that may be different to your grid. Start with "QUESTION:" and do not add anything else.} \\
            }
        }
    }
    & \\ \\

    \theutterance \stepcounter{utterance}  
    & & & \multicolumn{4}{p{0.6\linewidth}}{
        \cellcolor[rgb]{0.9,0.9,0.9}{
            \makecell[{{p{\linewidth}}}]{
                \texttt{\tiny{[GM$\langle$P2]}}
                \texttt{DESCRIPTION: A 5x5 grid with all cells filled with X except for the third row, which has blank spaces (▢) in the first, second, fourth, and fifth positions, leaving only the third column as X in that row. This creates a vertical line of X's in the third column and a horizontal line in the third row, but with the center cell (3,3) being X, forming a "T" shape.} \\
\\ 
\texttt{QUESTION: In your grid, is the center cell (3,3) an X?} \\
            }
        }
    }
     \\ \\

    \theutterance \stepcounter{utterance}  
    & & & \multicolumn{2}{p{0.3\linewidth}}{
        \cellcolor[rgb]{0.9,0.9,0.9}{
            \makecell[{{p{\linewidth}}}]{
                \texttt{\tiny{[GM$|$GM]}}
                \texttt{abort, first word: DESCRIPTION:, but expected QUESTION:} \\
            }
        }
    }
    & & \\ \\

\end{supertabular}
}

\end{document}
